%%%%%%%%%%%%%%%%%%%%%%%%%%%%%%%%%%%%%%
% LaTeX poster template
% Created by Nathaniel Johnston
% August 2009
% http://www.nathanieljohnston.com/2009/08/latex-poster-template/
%%%%%%%%%%%%%%%%%%%%%%%%%%%%%%%%%%%%%%

\documentclass[final]{beamer}
\usepackage[scale=1.24]{beamerposter}
\usepackage{graphicx}			% allows us to import images

%-----------------------------------------------------------
% Define the column width and poster size
% To set effective sepwid, onecolwid and twocolwid values, first choose how many columns you want and how much separation you want between columns
% The separation I chose is 0.024 and I want 4 columns
% Then set onecolwid to be (1-(4+1)*0.024)/4 = 0.22
% Set twocolwid to be 2*onecolwid + sepwid = 0.464
%-----------------------------------------------------------

\newlength{\sepwid}
\newlength{\onecolwid}
\newlength{\twocolwid}
\newlength{\threecolwid}
\setlength{\paperwidth}{48in}
\setlength{\paperheight}{36in}
\setlength{\sepwid}{0.024\paperwidth}
\setlength{\onecolwid}{0.22\paperwidth}
\setlength{\twocolwid}{0.464\paperwidth}
\setlength{\threecolwid}{0.708\paperwidth}
\setlength{\topmargin}{-0.5in}
\usetheme{confposter}
\usepackage{exscale}

%-----------------------------------------------------------
% The next part fixes a problem with figure numbering. Thanks Nishan!
% When including a figure in your poster, be sure that the commands are typed in the following order:
% \begin{figure}
% \includegraphics[...]{...}
% \caption{...}
% \end{figure}
% That is, put the \caption after the \includegraphics
%-----------------------------------------------------------

\usecaptiontemplate{
\small
\structure{\insertcaptionname~\insertcaptionnumber:}
\insertcaption}

%-----------------------------------------------------------
% Define colours (see beamerthemeconfposter.sty to change these colour definitions)
%-----------------------------------------------------------

\setbeamercolor{block title}{fg=ngreen,bg=white}
\setbeamercolor{block body}{fg=black,bg=white}
\setbeamercolor{block alerted title}{fg=white,bg=dblue!70}
\setbeamercolor{block alerted body}{fg=black,bg=dblue!10}

%-----------------------------------------------------------
% Name and authors of poster/paper/research
%-----------------------------------------------------------

\title{CODE-RADE}
\author{Sean Murray, Bruce Becker}
\institute{CSIR Meraka, University of Cape Town}

%-----------------------------------------------------------
% Start the poster itself
%-----------------------------------------------------------

\begin{document}
\begin{frame}[t]
  \begin{columns}[t]												% the [t] option aligns the column's content at the top
    \begin{column}{\sepwid}\end{column}			% empty spacer column
    \begin{column}{\onecolwid}
        \begin{block}{The problem CODE-RADE solves}
             \center{The researcher is faced with a paradox.}\\
             \center\includegraphics[scale=.5]{img/legoparadox.png}\\
             There is a huge amount of computing resources available, and a rich, well maintained, interoperable infrstructure of services... \\
             \center\textbf BUT \\
             Getting applictions onto it requires far too much effort.
        \end{block}         
            \begin{block}{The problem solved by CODE-RADE}
             \begin{itemize}
                 \item Lowers barrier to entry to grid or cloud infrastructure, or a single HPC site.
                 \item Application expert can prove to resource provider that the application will run on the sites execution environment
                 \item Easily manages the lifecycle of applications across multiple versions, architectures and configurations
                 \item Ensures that once an application is certified, its available on as many sites as possible.
                 \item Promotes collaboration between researcher, research software engineer, and infrastructure provider.
             \end{itemize}
            \end{block}
      \vskip2ex
         \begin{block}{CODE-RADE is }
             \begin{itemize}
                     \item Cross platform
                         \begin{itemize}
                                 \item Build and test artifacts for an arbitrary set of targets.
                                 \item Promote diversity in computing platforms.
                                 \item Ensure proper optimisations and application portability
                         \end{itemize}
                     \item Atomic
                         \begin{itemize}
                                 \item Fine grained control over dependencies, versions, and targets
                                 \item Relevant action taken on each event of CODE-RADE cycle.
                         \end{itemize}
                     \item Community
                         \begin{itemize}
                                 \item No restriction on applications that can be integrated.
                                 \item Anyone can contribute applications, resource, code review, etc.
                         \end{itemize}
                     \item Automated
                         \begin{itemize}
                                 \item Heavy reliance on automated agents to reduce bias, lead time.
                                 \item \textbf User-driven
                         \end{itemize}
             \end{itemize}

         \end{block}
    \end{column}
    \begin{column}{\sepwid}\end{column}			% empty spacer column
    \begin{column}{\onecolwid}
         \begin{block}{Achitecture}
             \begin{itemize}
                     \item Components
                           \begin{itemize}
                               \item Version Control (Github)
                               \item Continuous Integration (Jenkins)
                               \item Automated Delivery (CVMFS)
                           \end{itemize}
                     \item Workflow 
                         \begin{itemize}
                                 \item Build : User/expert-defined means to produce executable applications or libraries
                                 \item Test : Infrastructure operator-defined targets means to ensure viability
                                 \item Deliver : Ensure once tests pass, the application is available.
                         \end{itemize}
             \end{itemize}

         \end{block}
         \begin{block}{Target Matrix}
             \begin{itemize}
                 \item One can add as many target configurations to the build matrix as can be reliably simulated.
                 \item Currently we have defined three axes :
                 \begin{itemize}
                    \item Architecture x86\_64, ARM, etc.
                    \item Operating System RHEL Ubuntu, SLES, etc.
                    \item site specific - an axis encapsulating the special nature of a site, its interconnects, compiler, local libraries.
                 \end{itemize} 
             \end{itemize}
        \end{block}
         \begin{block}{Technology}
             \begin{itemize}
                 \item docker
                 \item quay.io
                 \item persistant storage
                 \item releases get a doi, authors are submitters and PR
             \end{itemize}
         \end{block}
    \end{column}
%     \end{columns}
    \begin{column}{\onecolwid}
      \begin{block}{A Standard Block}
        The standard block environment looks like this. It has justified text and a green title with an underline. You can create one like so:
        \begin{semiverbatim}
          {\color{red}\\begin}\{{\color{blue}block}\}\{Title\}\newline
          \hskip1ex.......\newline
          {\color{red}\\end}\{{\color{blue}block}\}
        \end{semiverbatim}
      \end{block}
      \vskip2ex
      \begin{alertblock}{An ``Alert'' Block}
        The ``alert'' block environment looks like this. It also has justified text, but it has a border and a light background to make it stand out. You can create one like so:
        \begin{semiverbatim}
          {\color{red}\\begin}\{{\color{blue}alertblock}\}\{Title\}\newline
          \hskip1ex.......\newline
          {\color{red}\\end}\{{\color{blue}alertblock}\}
        \end{semiverbatim}
      \end{alertblock}
    \end{column}

    \begin{column}{\sepwid}\end{column}			% empty spacer column
    \begin{column}{\twocolwid}					  % create a three-column-wide column and then we will split it up later
      \begin{block}{Altering Column Spans}
        You can make columns that span multiple other columns relatively easily. Lengths are defined in the template that make columns look normal-ish if you want to use a four-column layout like this poster. If you want to use a different number of columns, you will have to modify those lengths accordingly at the top of the poster.tex file.
        
        In particular, near the top of the TeX file you will see lines that look like:
        \begin{semiverbatim}
          \hskip1ex\\setlength\{\\sepwid\}\{0.024\\paperwidth\}
          
          \hskip1ex\\setlength\{\\onecolwid\}\{0.22\\paperwidth\}
          
          \hskip1ex\\setlength\{\\twocolwid\}\{0.464\\paperwidth\}
          
          \hskip1ex\\setlength\{\\threecolwid\}\{0.708\\paperwidth\}
        \end{semiverbatim}
        
        Set ``sepwid'' to be some small length somewhere near 0.025 (this is the space between columns). Then if $n$ is the number of columns you want, you should set
        \begin{align*}
          \text{onecolwid} & = \frac{1}{n}(1-(n+1)\times\text{sepwid}), \\
          \text{twocolwid} & = 2\times\text{onecolwid} + \text{sepwid}, \\
          \text{threecolwid} & = 3\times\text{onecolwid} + 2\times\text{sepwid}.
        \end{align*}
      \end{block}
      \begin{columns}[t,totalwidth=\threecolwid]	% split up that three-column-wide column
        \begin{column}{\onecolwid}
          \setbeamercolor{block title}{fg=red,bg=white}%frame color
          \setbeamercolor{block body}{fg=black,bg=white}%body color
          \begin{block}{Block Colours}
            For the standard blocks there are two colours; one for the title and one for the block body:\\
            \begin{semiverbatim}
              {\color{red}\\setbeamercolor}\{block title\}\newline \{fg=red,bg=white\}
            \end{semiverbatim}
            \begin{semiverbatim}
              {\color{red}\\setbeamercolor}\{block  body\}\newline \{fg=black,bg=white\}
            \end{semiverbatim}
            The \emph{fg} colour sets the text colour and \emph{bg} sets the background colour.
            For the normal blocks it makes no sense to use a background colour other than white. You \emph{can} change it, but it will look weird!
          \end{block}
        \end{column}
        \begin{column}{\onecolwid}
          \setbeamercolor{block alerted title}{fg=black,bg=norange}	% frame color
          \setbeamercolor{block alerted body}{fg=black,bg=white}		% body color
          \begin{alertblock}{Alert Block Colours}
            You can similarly modify the colours for alert blocks (but try not to overdo it):\\
            \begin{semiverbatim}
              {\color{red}\\setbeamercolor}\{block title\}\newline \{fg=black,bg=norange\}
            \end{semiverbatim}
            \begin{semiverbatim}
              {\color{red}\\setbeamercolor}\{block  body\}\newline \{fg=black,bg=white\}
            \end{semiverbatim}
          \end{alertblock}        
        \end{column}
        \begin{column}{\onecolwid}
          \begin{block}{References}
            Some references and a graphic to show you how it's done:
            
		        \small{\begin{thebibliography}{99}
		        \bibitem{KLPL06} D.~W. Kribs, R. Laflamme, D. Poulin, M. Lesosky, Quantum Inf. \& Comp. \textbf{6} (2006), 383-399.
		        \bibitem{zanardi97} P. Zanardi, M. Rasetti, Phys. Rev. Lett. \textbf{79},  3306 (1997).
		        \end{thebibliography}}
			      \vspace{0.75in}
			      \begin{center}
			        \includegraphics[width=5in]{canada.jpg}
			      \end{center}
		      \end{block}
        \end{column}
      \end{columns}
      \vskip2.5ex
    \end{column}
  \begin{column}{\sepwid}\end{column}			% empty spacer column
 \end{columns}
\end{frame}
\end{document}
